\section*{Вступ}

\subsection*{Актуальність теми}

Сучасний академічний ландшафт перебуває у стані глибинних трансформацій, зумовлених конвергенцією цифрових технологій, зміною моделей фінансування досліджень та новими викликами глобальної наукової комунікації. За даними OpenAlex, світова база наукових публікацій перевищила 250 мільйонів праць, а щорічний приріст становить понад 5 мільйонів нових робіт \cite{OpenAlex2024Docs}. Цей експоненційний зріст обсягу наукової інформації створює безпрецедентний тиск на адміністративні структури закладів вищої освіти (ЗВО), змушуючи їх переосмислити традиційні підходи до управління дослідницькою діяльністю.

Традиційна парадигма адміністрування науки, яка десятиліттями спиралася на ретроспективний аналіз даних та реактивну звітність, досягла межі своєї ефективності. Дослідження OCLC Research засвідчує, що у більшості університетів управління дослідницькою інформацією відбувається децентралізовано на рівні факультетів, що призводить до фрагментації даних, дублювання зусиль та неможливості стратегічного планування \cite{Bryant2017RIMDefining}. В умовах глобальної конкуренції за грантові ресурси та талановитих дослідників адміністративні підрозділи ЗВО змушені трансформуватися від «реєстраторів досягнень» до «архітекторів наукового успіху».

Ключовим драйвером цієї трансформації виступає генеративний штучний інтелект (Generative AI, GenAI). На відміну від попередніх поколінь цифрових інструментів, що фокусувалися на зберіганні та класифікації даних, GenAI відкриває можливості для створення нового контенту, прогнозування трендів, автономного виконання складних аналітичних завдань та моделювання сценаріїв розвитку наукових екосистем \cite{Wang2024SurveyLLMAgents,Xi2023RisePotentialAgents}. Поява протоколу Model Context Protocol (MCP) у листопаді 2024 року та його передача до Agentic AI Foundation (Linux Foundation) у грудні 2025 року з підтримкою Google, Microsoft, AWS та Cloudflare створюють стандартизовану основу для інтеграції ШІ-інструментів з корпоративними інформаційними системами \cite{Anthropic2024MCP}.

Для української наукової спільноти проблема цифрової трансформації адміністрування науки набуває особливої гостроти в контексті євроінтеграційних процесів, необхідності відповідності стандартам OpenAIRE та European Open Science Cloud (EOSC), а також обмежених кадрових та фінансових ресурсів \cite{OpenAIRE2023GuidelinesCRIS}. Водночас партнерство Міністерства цифрової трансформації України з University College London у сфері AI governance, інвестиції World Bank у цифрову інфраструктуру та поява нових фахових видань з фокусом на цифрових практиках публічного управління створюють сприятливе середовище для впровадження інноваційних рішень.

\subsection*{Проблема дослідження}

Незважаючи на значний інтерес дослідників до застосування ШІ у вищій освіті, існуючі наукові розвідки переважно зосереджуються на використанні GenAI в навчальному процесі --- для генерації контенту, оцінювання студентів, персоналізації навчання [30, 32]. Адміністративна сфера управління наукою залишається значною мірою поза фокусом досліджень, хоча саме тут накопичений колосальний масив рутинних операцій, які можуть бути автоматизовані.

Ключові проблеми, що потребують вирішення:

\textbf{По-перше, фрагментація даних.} Наукометричні показники дослідників зберігаються в ізольованих системах --- ORCID, Scopus, Web of Science, Google Scholar, OpenAlex, локальні бази ЗВО. Відсутність єдиного «золотого запису» (Golden Record) призводить до систематичного дублювання публікацій (за нашими оцінками, 15--25\% записів у типових звітах є дублікатами), неузгодженості метаданих та неможливості комплексного аналізу.

\textbf{По-друге, трудомісткість звітності.} Підготовка річних звітів про наукову діяльність кафедри традиційно вимагає 8--12 годин ручної праці --- копіювання даних з різних джерел, верифікація, форматування відповідно до вимог МОН України. Ця рутинна робота поглинає ресурси, які могли б бути спрямовані на стратегічні завдання.

\textbf{По-третє, реактивність.} Звіти формуються post factum, унеможливлюючи оперативне реагування на тенденції (наприклад, зниження публікаційної активності окремих дослідників чи кафедр) та прогнозування. Управління залишається ретроспективним, орієнтованим на compliance замість розвитку.

\textbf{По-четверте, відсутність науково обґрунтованих механізмів інтеграції GenAI} в адміністративні процеси ЗВО. Існуючі рішення або є закритими комерційними продуктами (Pure, Converis), або не враховують специфіку автономних ШІ-агентів, здатних самостійно виконувати послідовності дій для досягнення мети.

\subsection*{Мета та завдання дослідження}

\textbf{Мета дослідження} --- розробка та апробація концептуальної моделі та механізмів інтеграції генеративного штучного інтелекту для переходу від реактивної системи бібліометричної звітності до проактивного управління науковою діяльністю ЗВО.

Для досягнення мети визначено такі \textbf{завдання}:
\begin{enumerate}[leftmargin=*,itemsep=0.2em]
\item Здійснити концептуальний аналіз парадигм реактивного та проактивного управління в контексті адміністрування науки, визначити критерії та індикатори трансформації.
\item Систематизувати технічні можливості GenAI для автоматизації наукового адміністрування, включаючи обробку бібліометричних даних, предиктивну аналітику та документообіг.
\item Розробити механізми технічної інтеграції GenAI з інформаційними системами ЗВО на основі протоколу MCP та підходу «Golden Record».
\item Провести практичну апробацію запропонованих механізмів на базі системи ResearchIQ та оцінити кількісні та якісні результати.
\item Проаналізувати світовий досвід впровадження ШІ в управлінні дослідженнями та визначити best practices.
\item Ідентифікувати ризики, обмеження та етичні аспекти використання автономних ШІ-агентів у публічному адмініструванні.
\item Сформулювати рекомендації та дорожню карту впровадження GenAI для українських ЗВО.
\end{enumerate}

\subsection*{Об'єкт та предмет дослідження}

\textbf{Об'єкт дослідження} --- процеси адміністрування наукової діяльності в закладах вищої освіти.

\textbf{Предмет дослідження} --- механізми інтеграції генеративного штучного інтелекту для цифрової трансформації бібліометричної звітності та переходу до проактивного управління.

\subsection*{Наукова новизна}

Наукова новизна одержаних результатів полягає в наступному.

\textbf{Вперше:}
\begin{itemize}[leftmargin=*,itemsep=0.2em]
\item обґрунтовано та реалізовано концепцію \textbf{MCP Function Calling Agent} для автономної обробки наукометричних даних, де ШІ-агент самостійно визначає послідовність API-запитів до зовнішніх баз (ORCID, OpenAlex) на основі декларативного опису доступних функцій (tools);
\item запропоновано трирівневу архітектуру інтеграції GenAI з CRIS-системами (Data Layer $\to$ Agent Layer $\to$ Interface Layer), що забезпечує масштабованість та відповідність стандартам CERIF/OpenAIRE.
\end{itemize}

\textbf{Удосконалено:}
\begin{itemize}[leftmargin=*,itemsep=0.2em]
\item алгоритм мультиджерельного злиття даних \textbf{«Golden Record»} шляхом додавання механізмів нормалізації DOI, двоетапного матчингу (DOI + Title/Year) та атрибутної агрегації з трекінгом походження даних (provenance);
\item методичні підходи до оцінки ефективності цифрової трансформації адміністративних процесів через систему кількісних (час, точність, охоплення) та якісних (аналітична цінність, user experience) індикаторів.
\end{itemize}

\textbf{Набули подальшого розвитку:}
\begin{itemize}[leftmargin=*,itemsep=0.2em]
\item теоретичні положення концепції проактивного управління (proactive management) в контексті публічного адміністрування науки;
\item рекомендації щодо етичного використання автономних ШІ-агентів у публічному секторі з урахуванням вимог EU AI Act та GDPR.
\end{itemize}

\subsection*{Структура статті}

Стаття структурована відповідно до логіки Design Science Research. У першому розділі викладено теоретико-методологічні засади дослідження, включаючи концептуальний апарат та обґрунтування методології DSR. Другий розділ присвячено аналізу проблемного поля --- кризи реактивної моделі бібліометричної звітності. Третій розділ систематизує технічні можливості GenAI для автоматизації наукового адміністрування. Четвертий розділ детально розглядає механізми технічної інтеграції (CRIS, MCP, Golden Record). П'ятий розділ представляє результати практичної апробації системи ResearchIQ. Шостий розділ аналізує світовий досвід впровадження. Сьомий розділ присвячено ризикам, обмеженням та етичним аспектам. Висновки узагальнюють результати та визначають перспективи подальших досліджень.